\documentclass[11pt,a4paper]{article}
\usepackage[T1]{fontenc}
\usepackage[utf8]{inputenc}
\usepackage{amsmath,amssymb,amsthm}
\usepackage[margin=1in]{geometry}
\usepackage[colorlinks=true,linkcolor=blue,urlcolor=blue]{hyperref}
\theoremstyle{plain}
\newtheorem{theorem}{Theorem}
\newtheorem{lemma}[theorem]{Lemma}
\newtheorem{proposition}[theorem]{Proposition}
\newtheorem{corollary}[theorem]{Corollary}
\theoremstyle{definition}
\newtheorem{remark}[theorem]{Remark}
% A quoted conjecture can be a single hundred-term formula whose terms are thirty-digit
% numbers. TeX cannot hyphenate those, so without a little stretch every such quote runs off
% the page; the linked-a-file papers are all of that shape.
\setlength{\emergencystretch}{3em}

\title{The conjectured generating function of OEIS A065261,\\ \large from the recurrence the entry
states}
\author{Adrian Perez Fontelles\\ \small Independent researcher}
\date{15 September 2026}
\begin{document}
\maketitle

\begin{abstract}
OEIS A065261 states a linear recurrence as a fact and carries a rational generating function as an
empirical conjecture. The conjecture follows from the recurrence by polynomial algebra alone:
the recurrence fixes the denominator, the published terms fix the numerator exactly, and the
generating function is then determined as a rational function rather than checked on finitely
many coefficients. No model of the sequence and no reading of its name is required.
\end{abstract}

\noindent\small 2020 Mathematics Subject Classification. 11B37, 05A15, 30B10.\normalsize

\section{What the entry says}

The first terms are $1, 1, 1, 2, 5, 3, 2, 4,\dots$, with offset $1$. The entry states, as a fact,

\begin{quote}\raggedright a(n) = 2*a(n-4)-a(n-8).\end{quote}

\noindent and separately records the conjecture

\begin{quote}\raggedright Empirical g.f.: x*(x\^{}5+3*x\^{}4+2*x\^{}3+x\^{}2+x+1) / ((x-1)\^{}2*(x+1)\^{}2*(x\^{}2+1)\^{}2). [\_Colin Barker\_, Feb 18 2013]\end{quote}

The stated line was taken from the entry's factual formulas only. A formula qualified by a
finite range -- "holds at least up to $n=1000$ but is not known to hold in general" -- is not a
fact and is not used, nor is anything inside a conjectural block.

\section{The argument}

Write $b_k=a(1+k)$ and
\[
  D(x)=1-\sum_{i} c_i x^{i}=x^{8} - 2 x^{4} + 1 ,
\]
the polynomial of the stated recurrence, of degree $8$. Let $B(x)=\sum_{k\ge0}b_kx^k$.

The recurrence says $b_k=\sum_i c_ib_{k-i}$ for every $k\ge 8$. The coefficient of $x^k$ in
$D(x)B(x)$ is exactly $b_k-\sum_i c_ib_{k-i}$, so it vanishes for every $k\ge 8$ and
\[
  D(x)B(x)=P(x)=\sum_{k<8}\Bigl(b_k-\sum_i c_ib_{k-i}\Bigr)x^k
         =x^{6} + 3 x^{5} + 2 x^{4} + x^{3} + x^{2} + x ,
\]
a polynomial whose every coefficient is fixed by finitely many published terms. Therefore
\[
  B(x)=\frac{P(x)}{D(x)}=\frac{x \left(x^{5} + 3 x^{4} + 2 x^{3} + x^{2} + x + 1\right)}{\left(x - 1\right)^{2} \left(x + 1\right)^{2} \left(x^{2} + 1\right)^{2}} ,
\]
which is the conjectured expression.

\begin{theorem}
The conjectured generating function of OEIS A065261 is correct, given the recurrence the entry
states.
\end{theorem}

\begin{proof}
Immediate from the displayed identity: $D\,B=P$ holds coefficientwise, $D(0)=1\ne0$ so $D$ is
invertible in the ring of formal power series, and $B=P/D$ follows. This is an identity of
rational functions, not a comparison of the $76$ terms the entry publishes -- those
enter only in fixing the $8$ coefficients of $P$.
\end{proof}

\end{document}
